%% parts/p06/cardinal-arithmetic.tex --- Phase-4 P6 cardinal arithmetic chapter
%% Expands the cardinal material migrated from v1.1 lines 10559--10654.
%% Depth: D-P7-002/003/004 and N-a-035; apparatus: A-061.
%% N-a-036 remains an explicit could-row stub at its S ceiling; its PCF
%% continuation is excluded by N-x-005.
\chapter{Cardinal Arithmetic}
\label{sec:p06-cardinal-arithmetic}
\chaptermark{Cardinal Arithmetic}

Throughout this chapter, cardinal arithmetic means arithmetic of
ordinal cardinals in an ambient theory containing the Axiom of Choice, unless
a statement explicitly says otherwise.

\begin{hbdefinition}{def:p06-cardinal-arithmetic}{Cardinal sum, product, and exponentiation}
For cardinals \(\kappa,\lambda\), define
\[
\begin{aligned}
  \kappa+\lambda
    &\defeq\abs{(\kappa\times\{0\})\cup(\lambda\times\{1\})},\\
  \kappa\cdot\lambda
    &\defeq\abs{\kappa\times\lambda},\\
  \kappa^\lambda
    &\defeq\abs{{}^{\lambda}\kappa},
\end{aligned}
\]
where \({}^{\lambda}\kappa\) is the set of functions from \(\lambda\) to
\(\kappa\).  For a set-indexed family \(\langle\kappa_i:i\in I\rangle\),
\(\sum_{i\in I}\kappa_i\) is the cardinal of the disjoint union and
\(\prod_{i\in I}\kappa_i\) is the cardinal of the Cartesian product.
These operations act on cardinals, not on arbitrary representatives; their
set constructions distinguish them from ordinal addition, multiplication,
and exponentiation.
\index{cardinal arithmetic}
\index{cardinal addition}
\index{cardinal multiplication}
\index{cardinal exponentiation}
\end{hbdefinition}

\begin{hbtheorem}{thm:p06-cardinal-arithmetic-laws}{cited=Jech2003}{Basic laws of cardinal arithmetic}
Assume the Axiom of Choice.  Cardinal addition and multiplication are
commutative and associative, multiplication distributes over addition, and
for cardinals \(\kappa,\lambda,\mu\) with the displayed expressions defined,
\[
  \kappa^{\lambda+\mu}=\kappa^\lambda\cdot\kappa^\mu,
  \qquad
  (\kappa\cdot\lambda)^\mu=\kappa^\mu\cdot\lambda^\mu,
  \qquad
  (\kappa^\lambda)^\mu=\kappa^{\lambda\cdot\mu}.
\]
If \(\kappa\) and \(\lambda\) are infinite, then
\[
  \kappa+\lambda=\kappa\cdot\lambda=\max(\kappa,\lambda).
\]
Moreover \(2^\kappa=\abs{\mathcal P(\kappa)}\) and
\(\kappa<2^\kappa\) for every cardinal \(\kappa\).
\end{hbtheorem}

\begin{hbdefinition}{def:p06-aleph-beth-continuum}{Alephs, beths, CH, and GCH}
Define the aleph hierarchy by
\[
  \aleph_0\defeq\omega,\qquad
  \aleph_{\alpha+1}\defeq(\aleph_\alpha)^+,
  \qquad
  \aleph_\lambda\defeq\sup_{\beta<\lambda}\aleph_\beta
\]
at a limit ordinal \(\lambda\), where \(\kappa^+\) is the least cardinal
strictly greater than \(\kappa\).  Define the beth hierarchy by
\[
  \beth_0\defeq\aleph_0,\qquad
  \beth_{\alpha+1}\defeq2^{\beth_\alpha},
  \qquad
  \beth_\lambda\defeq\sup_{\beta<\lambda}\beth_\beta.
\]
The \emph{Continuum Hypothesis} and \emph{Generalized Continuum Hypothesis}
are respectively
\[
  \mathsf{CH}\defeq2^{\aleph_0}=\aleph_1,
  \qquad
  \mathsf{GCH}\defeq
  \forall\alpha\,\bigl(2^{\aleph_\alpha}=\aleph_{\alpha+1}\bigr).
\]
The successor notation, aleph recursion, and beth recursion name cardinal
operations and hierarchies; CH and GCH are statements about them, not
definitions of the continuum function.
\index{aleph hierarchy}
\index{beth hierarchy}
\index{successor cardinal}
\index{Continuum Hypothesis}
\index{Generalized Continuum Hypothesis}
\end{hbdefinition}

\begin{hbdefinition}{def:p06-cofinality-regular-singular}{Cofinality and regular or singular cardinals}
A subset \(C\) of an ordinal \(\alpha\) is \emph{cofinal} in \(\alpha\) iff
for every \(\beta<\alpha\) some \(\gamma\in C\) satisfies
\(\beta\leq\gamma\).  The \emph{cofinality} \(\operatorname{cf}(\alpha)\)
is the least ordinal that is the order type of a cofinal subset of \(\alpha\),
equivalently the least ordinal admitting a cofinal map into \(\alpha\).
An infinite cardinal \(\kappa\) is \emph{regular} iff
\(\operatorname{cf}(\kappa)=\kappa\), and \emph{singular} iff
\(\operatorname{cf}(\kappa)<\kappa\).  Cofinality measures how many earlier
stages are required to approach an ordinal; it is not its cardinality.
\index{cofinal subset}
\index{cofinality}
\index{regular cardinal}
\index{singular cardinal}
\end{hbdefinition}

\begin{hbtheorem}{thm:p06-cofinality-facts}{cited=Jech2003}{Basic cofinality facts}
Assume the Axiom of Choice.  For every nonzero limit ordinal \(\alpha\),
\(\operatorname{cf}(\alpha)\) is an infinite regular cardinal and
\(\operatorname{cf}(\alpha)\leq\abs\alpha\).  Every infinite successor
cardinal is regular.  A cardinal \(\kappa\) is singular exactly when it is the
union of fewer than \(\kappa\) many sets, each of cardinality below
\(\kappa\).
\end{hbtheorem}

\begin{hbtheorem}{thm:p06-konig-cardinal}{cited=Konig1905}{K\"onig's theorem}
Assume the Axiom of Choice.  Let \(I\) be a set and let
\(\langle\kappa_i:i\in I\rangle\) and
\(\langle\lambda_i:i\in I\rangle\) be families of cardinals such that
\(\kappa_i<\lambda_i\) for every \(i\in I\).  Then
\[
  \sum_{i\in I}\kappa_i
  <
  \prod_{i\in I}\lambda_i.
\]
In particular, for every infinite cardinal \(\kappa\),
\[
  \operatorname{cf}(2^\kappa)>\kappa
  \qquad\text{and}\qquad
  \kappa^{\operatorname{cf}(\kappa)}>\kappa.
\]
\end{hbtheorem}

\begingroup\small
\begin{longtable}{@{}p{.22\textwidth}p{.37\textwidth}p{.32\textwidth}@{}}
\caption{Cardinal-arithmetic register (apparatus A-061).}
\label{tab:p06-cardinal-arithmetic}\\
\toprule
Entry & Formula or classification & Hypothesis and use \\
\midrule
\endfirsthead
\toprule
Entry & Formula or classification & Hypothesis and use \\
\midrule
\endhead
Finite sums and products & Definitions in
Definition~\ref{def:p06-cardinal-arithmetic} & Set constructions; ordinal
cardinal assignment uses the local Choice standing. \\
Infinite collapse & \(\kappa+\lambda=\kappa\cdot\lambda=
\max(\kappa,\lambda)\) for infinite \(\kappa,\lambda\) & Axiom of Choice;
supplies the finite-coding bounds used in Part~1. \\
Exponentiation & \(\kappa^\lambda=\abs{{}^{\lambda}\kappa}\),
\(2^\kappa=\abs{\mathcal P(\kappa)}\), and the three exponent laws above & Axiom
of Choice for ordinal-cardinal assignment; Cantor gives
\(\kappa<2^\kappa\). \\
Aleph hierarchy & \(\aleph_0=\omega\), successor by \(\kappa^+\), limit by
supremum & Enumerates the infinite ordinal cardinals under Choice. \\
Beth hierarchy & \(\beth_0=\aleph_0\), successor by power set, limit by
supremum & Tracks iterated power-set cardinalities. \\
Cofinality & Least order type of a cofinal subset & Separates regular
\(\operatorname{cf}(\kappa)=\kappa\) from singular
\(\operatorname{cf}(\kappa)<\kappa\) cardinals. \\
K\"onig's theorem & \(\sum_i\kappa_i<\prod_i\lambda_i\) when
\(\kappa_i<\lambda_i\) & Forces
\(\operatorname{cf}(2^\kappa)>\kappa\); constrains every continuum function. \\
Hausdorff successor recursion &
\((\kappa^+)^\lambda=\kappa^\lambda\cdot\kappa^+\) for infinite
\(\kappa\) and nonzero \(\lambda\) & Reduces a successor base to its
predecessor; the nonzero-exponent hypothesis is essential. \\
CH and GCH & \(2^{\aleph_0}=\aleph_1\); respectively
\(2^{\aleph_\alpha}=\aleph_{\alpha+1}\) for every \(\alpha\) & Independent
of \(\mathsf{ZFC}\); the forcing chapter records the calibrated status. \\
Easton regular-cardinal freedom & Prescribed values of \(2^\kappa\) on
infinite regular cardinals, subject to monotonicity and the cofinality bound &
Theorem~\ref{thm:p06-easton} states the exact cited hypotheses over a
\(\mathsf{ZFC}+\mathsf{GCH}\) ground; no prescription at singular cardinals is
asserted. \\
\bottomrule
\end{longtable}
\endgroup

%% Source exercises 10656--10673 are deliberately omitted because the
%% handbook has no exercise layer; their ordinal, rank, and collapse topics
%% are entries in the earlier P6 axioms-and-hierarchy chapter.
\begin{hbtheorem}{thm:p06-hausdorff-recursion}{cited=Hausdorff1904}{Hausdorff's successor recursion}
Assume the Axiom of Choice.  For every infinite cardinal \(\kappa\) and every
nonzero cardinal \(\lambda\),
\[
  (\kappa^+)^\lambda
  =\kappa^\lambda\cdot\kappa^+
  =\max(\kappa^\lambda,\kappa^+).
\]
Equivalently, for ordinals \(\alpha,\beta\),
\[
  \aleph_{\alpha+1}^{\aleph_\beta}
  =\aleph_\alpha^{\aleph_\beta}\cdot\aleph_{\alpha+1}.
\]
The restriction \(\lambda\neq0\) cannot be dropped: the left side would be
\(1\) while the displayed product would be \(\kappa^+\).  Repeated use of
this identity is the Hausdorff recursion for exponentiation with a successor
aleph base \cite{Jech2003}.
\end{hbtheorem}

\begin{hbtheorem}{thm:p06-easton}{cited=Easton1970}{Easton's theorem on regular cardinals}
Let \(F\) be a definable class function from the infinite regular cardinals to
cardinals such that
\(\kappa<\lambda\) implies \(F(\kappa)\leq F(\lambda)\), and
\(\operatorname{cf}(F(\kappa))>\kappa\) for every regular \(\kappa\).
Starting from a ground model of \(\mathsf{ZFC}+\mathsf{GCH}\) carrying such an
\(F\), Easton's forcing yields a cofinality-preserving generic extension in
which \(2^\kappa=F(\kappa)\) at every infinite regular cardinal \(\kappa\).
Thus monotonicity and K\"onig's cofinality condition are, relative to the
existence of the stated ground model, the only restrictions on the continuum
function at regular cardinals.  The theorem does not prescribe exponentiation
at singular cardinals and is a relative-consistency result, not a theorem that
\(\mathsf{ZFC}\) chooses one such function in the ground universe.
\end{hbtheorem}


\hbstub{N-a-036 remains owed at its S ceiling: add Silver's theorem statement.  PCF theory beyond that statement remains excluded by N-x-005.}
